
\section{Benchmark and Experimental Design}
\label{sec:method}

\subsection{Fresh formal benchmark}

We construct 180 fresh scenarios: six task families with 30 samples each, distributed
across 12 domains. Each sample contains a machine-readable \texttt{world\_spec}, source
records with timestamps and provenance, a rendered decision question, distractor evidence,
explicit required facts, permissible action families, and a safe fallback when obligations
are uncovered.

\begin{table}[t]
\centering
\small
\setlength{\tabcolsep}{4pt}
\begin{tabular}{p{0.08\linewidth}p{0.23\linewidth}p{0.31\linewidth}p{0.30\linewidth}}
\toprule
Family & Obligation & C1 visible-gap cue & Silent-failure route \\
\midrule
F1 & Constraint / feasibility & Requirement remains while feasibility support is missing & Top-$k$ can remove cue-bearing context together with support \\
F2 & Authority / approval & Authority requirement remains without current grant & Top-$k$ can surface action support without the requirement \\
F3 & Temporal validity & Stale marker remains without current replacement & Not observed under frozen retrievers \\
F4 & Conflict resolution & Conflicting claims remain without resolution & Top-$k$ can remove one side or the resolution structure \\
F5 & Outcome / revision & Prior conditional decision remains without change evidence & Not observed under frozen retrievers \\
F6 & Override / invalidation & None when later override is removed & Older instruction remains locally actionable \\
\bottomrule
\end{tabular}
\caption{Task families and policy-relative gap structure. The last column describes a
possible silent surface under the frozen policy; it is not a universal property of the
family.}
\label{tab:families}
\end{table}

The benchmark is independent of the earlier TMLR validation benchmark and sealed Holdback;
neither is reused as confirmatory evidence here.

\subsection{Formal validation and freeze}

Every formal sample must pass two separately implemented validators:
(1) a declarative evaluator reading the frozen \texttt{world\_spec.decision\_contract};
and (2) a family-specific evaluator that independently derives required fact types and
action families. They must agree exactly on required fact IDs, mandatory evidence IDs,
full-evidence action families, fallback action families, and coverage completeness.

Each sample also passes metamorphic controls:
FULL covers every required obligation; hiding the designated treatment target makes at
least one required obligation uncovered; removing a declared non-mandatory evidence item
preserves complete coverage; and world truth is unchanged across the paired treatment.

Before P2, the frozen dataset had zero schema failures, zero dual-validator disagreements,
zero renderer/world-spec consistency failures, zero dataset-audit errors, and zero
near-duplicate warnings.

\subsection{Frozen downstream policy}

P2 uses one deterministic policy frozen before outcomes. In C0--C4 it can read only the
decision question, visible evidence role, evidence currentness, candidate support,
candidate actions, and prior-state fields. It cannot read \texttt{world\_spec}, Gold,
condition labels, treatment targets, or construction metadata.

Its visible-gap rules are frozen as:
\begin{itemize}
  \item visible constraint requirement with no feasibility evidence $\rightarrow$ \textsc{Clarify};
  \item visible authority requirement with no current grant $\rightarrow$ \textsc{RequestApproval};
  \item visible stale state with no current state $\rightarrow$ \textsc{Defer};
  \item visible conflicting claims with no resolution $\rightarrow$ \textsc{Clarify};
  \item visible prior conditional decision with no outcome-change evidence $\rightarrow$ \textsc{Defer}.
\end{itemize}
Otherwise the policy chooses a decisive action from visible candidate-support evidence
using frozen role and currentness weights.

\subsection{Conditions}

\begin{table}[t]
\centering
\small
\setlength{\tabcolsep}{4pt}
\begin{tabular}{p{0.12\linewidth}p{0.31\linewidth}p{0.20\linewidth}p{0.27\linewidth}}
\toprule
Condition & Evidence surface & Extra coverage input & Purpose \\
\midrule
C0 FULL & all formal evidence & none & positive control \\
C1 HIDDEN & one designated mandatory item removed & none & paired causal intervention \\
C2 LEXICAL & IDF token overlap, top-$k=2$ & none & controlled lexical retrieval \\
C3 HASH-DENSE & 256-D signed char-trigram hash, cosine, top-$k=2$ & none & controlled dense-vector baseline \\
C4 HYBRID & RRF of C2+C3, constant 60, top-$k=2$ & none & fixed hybrid surface \\
C5 COVERAGE-AWARE & \textbf{identical C4 surface} & \texttt{coverage\_complete} only & mechanism test \\
\bottomrule
\end{tabular}
\caption{Frozen experimental conditions. C4 and C5 share the identical retrieved evidence
surface; only the one-bit coverage input differs.}
\label{tab:conditions}
\end{table}

C3 is a deterministic dense-vector baseline, not a pretrained semantic embedding model.
No retriever parameter was tuned after observing P2 outcomes.

\subsection{Endpoints and statistics}

The primary endpoint is the paired risk difference
$\mathrm{UDR}(C1)-\mathrm{UDR}(C0)$. The preregistered test is a two-sided exact McNemar
test with $\alpha=0.05$. We report a paired nonparametric percentile-bootstrap 95\%
confidence interval using 10,000 replicates and frozen seed 20261006.

Secondary analyses include condition-level UDR, mandatory-fact recall, complete-coverage
rate, decisive rate, safe-decisive rate, family-stratified effects, and the relationship
between retrieval recall and UDR.

Because the scenarios are procedurally constructed rather than sampled from a defined
real-world task population, these inferential statistics summarize the paired benchmark
intervention under the frozen protocol. They must not be interpreted as estimates of
real-world failure prevalence.
